% AI-written, AI cross-reviewed; no human sign-off; not a claim that RH is proved
% Atlas unit P037. Local manuscript v3, item 7.14, heading on printed PDF p. 81.
% PositivityTheory.QuaternionShaBounds.sSha
% paper2/inputs_one.tex:623–658. Verbatim excerpt; SHA-256 f5b249c9dd5bdc930965ebf6ef94f1428daa1ad2bc783610082e69f22f1583bf.
% Original TeX uses the macros and cross-reference keys of the source manuscript.


% Start of atlas edition annotation. The following block is outside the original excerpt.
\paragraph{Status of named inputs in this reading.}
\noindent\nolinkurl{PositivityTheory.QuaternionShaBounds.ClInj_holds} --- in-repo constant (\path{Development/AutomorphicArithmetic/AutomorphicVarieties/ModularityArgument/ClassCounting/RingEquivalences/OrderOfOneAddPrimePowZmodPrimePow.lean:265--330}); not compared.\par
% End of atlas edition annotation.

\begin{theorem}[From a uniform layer bound to the localization kernel]\label{pn:P037}
Let \(L\) be totally complex and \(\ell\) prime. Suppose \(\ell^c\)
is the largest \(\ell\)-power order of a root of unity in \(L\),
\(c\ge1\) and \(c\ge2\) at \(\ell=2\). Fix \(E\ge1\). If for every
allowed layer \(F/L\) the preceding character group, defined with
\(a=(1+\ell^c)^E\), is finite of order at most \(B\), then there
exists \(s\in\mathbb N\) such that
\[
 \operatorname{shaInv}(L,\ell,E)=s,\qquad \ell^s\le B .
\]
\end{theorem}\nodestatus{P037}
\begin{proof}
Restriction to the infinite layer is injective by the comparison in
Blueprint Z2.K1~\bploc{lem:P-r3-k1}. Represent a class by a cocycle \(z\). Its image
is finite. At each of the finitely many places above \(\ell\)
choose \(b_w\in A_E\) giving its local coboundary on the decomposition
group. Choose \(\nu\) such that \(\ell^\nu\) kills these elements
and the image of \(z\). Over \(F=L(\mu_{\ell^\nu})\) its restriction
is a homomorphism and is zero on the chosen local groups and their
conjugates. The unramified quotient imposes the remaining inertia
conditions. By unramified Artin reciprocity this character factors
through the ordinary class group; primes over \(\ell\) map to zero.
Under the identification of \(A_E[\ell^\nu]\) with \(\mathbb Z/\ell^\nu\),
the action of \(\gamma^{-1}\) gives the eigenvalue \(a\). Evaluations
on prime ideals determine the resulting character.

For any finite set of original classes one can choose a common
\(\nu\), taking the maximum of their thresholds. Distinct restrictions
at the infinite layer must then have distinct finite-level characters.
Every such set has size at most \(B\), so the localization kernel is
finite of order at most \(B\). A finite \(\mathbb Z_\ell\)-module of
length \(s\) has order \(\ell^s\). The twist, the common-level
argument, and this length comparison are the additions to reciprocity
in Blueprint Z2.K6~\bploc{lem:P-r3-k6}. See J.~S.~Milne, \emph{Class Field Theory},
Theorem~0.4 and Chapter~V. \cite{r21-cft}
\end{proof}
